\chapter{Introduction}
\lecture{1}{13 Oct. 08:00}{First Lecture}
\section{Theorem Environments}

\begin{definition}[Natural number]\label{def}
	The set of \vocab{natural numbers} is denoted \(\mathbb{N}\).
\end{definition}

\begin{lemma}[Useful lemma]\label{lma}
	Given the axioms of \hyperref[def]{natural numbers \(\mathbb{N}\)}, we have
	\[
		0 \neq 1.
	\]
\end{lemma}
\begin{proof}[An obvious proof]
	Obvious.
\end{proof}

\begin{proposition}[Useful proposition]\label{prop}
	From \autoref{lma}, we have
	\[
		0 < 1.
	\]
\end{proposition}

\begin{theorem}[Mass-energy equivalence]\label{thm}
	Given \autoref{prop}, we then have
	\[
		E = mc^2.
	\]
\end{theorem}
\begin{proof}
	The blank left for me is too small, hence we put the proof in \autoref{appendix}.
	\begin{subproof}[Proof of key step]
		This inner proof uses $\blacksquare$ to distinguish from the outer $\blacksquare$.
	\end{subproof}
\end{proof}

\begin{corollary}[Riemann hypothesis]\label{col}
	The real part of every nontrivial zero of the Riemann zeta function is \(\frac{1}{2}\).
\end{corollary}
\begin{proof}
	Left to the reader.\todo{DIY}
\end{proof}

\hrulebar

\section{Other Environments}

\begin{exercise}
	Prove that \(1 < 2\).
\end{exercise}
\begin{proof}[Solution]
	Iteratively apply \autoref{prop}.
\end{proof}

\begin{eg}[Interesting example]\label{eg}
	We note that \(1 \neq 2\)!
	\begin{note}[Important note]
		As a consequence, \(2 \neq 3\) also.
	\end{note}
\end{eg}

\begin{remark}
	See that \autoref{prop} is powerful. We also have a citation: \cite{newton1726philosophiae}.
\end{remark}

\begin{prev}
	We saw \autoref{lma} is helpful.
\end{prev}

\begin{claim}
	The center \(Z(G)\) is always a normal subgroup of \(G\).
\end{claim}
\begin{subproof}
	For any \(z \in Z(G)\) and \(g \in G\), we have \(gzg^{-1} = z \in Z(G)\).
\end{subproof}

\begin{intuition}
	\vocab{Key idea}: think of the center as the fixed points of conjugation.
\end{intuition}

\subsubsection{Internal links}
You can use \texttt{\textbackslash autoref} for internal references and \texttt{\textbackslash cite} for bibliography entries.

\section{Footnote Test}

\begin{theorem}
	Every group of prime order is cyclic.\footnote{This is a footnote inside a theorem box.}
\end{theorem}

\begin{eg}
	The integers \(\mathbb{Z}\) form an infinite cyclic group.\footnote{This is a footnote inside an example box.}
\end{eg}

\section{Commutative Diagrams}

\begin{eg}
	A commutative square drawn with \texttt{tikz-cd}:
	\[
		\begin{tikzcd}
			A \arrow[r, "f"] \arrow[d, "g"'] & B \arrow[d, "h"] \\
			C \arrow[r, "k"']                & D
		\end{tikzcd}
	\]
\end{eg}

\section{Vocab and Dividers}

Use \texttt{\textbackslash vocab\{term\}} to highlight a \vocab{key term} in blue bold.

Use \texttt{\textbackslash hrulebar} to insert a visual section divider:

\hrulebar

Content continues after the divider.
